% Nouveau EGA V (cote 25): archive cote 25, site page 15 (PDF page 16).
% typescript. Draft compared: none.
% AI-transcribed by Claude Opus 5.5 (claude-opus-5-5, adaptive thinking, effort high), Message Batches API, batch ega25b (msgbatch_01A7n5YuXubUPNK1zsepi8RG), 2026-11-13; not yet checked by eye.
\scanpage{25/15}
\mspage{12}

\uline{Corollaire} 2.17.   Soit Z l'ensemble des points de X en lesquels F
n'est pas séparable sur k (réf). Alors  $G_{\hand{\unsure{y|\eta}}}$ est séparable sur $k(\hand{\unsure{y|\eta}})$ sss
Z est fini, En particulier, si F est séparable sur k, $G_{\hand{\unsure{y|\eta}}}$ est séparable
sur $k(\hand{\unsure{y|\eta}})$.

\placedmargin{\unsure{bb pour|à mettre} dans l'énoncé 2.12.}\uline{Remarques} 2.18.  Le résultat clef  2.12. (et ses conséquences \unsure{2.13.|2.11.} à 2.17.) devient faux si on abandonne
l'hypothèse que f est une immer\unsure{z|s}ion, comme  on voit par exemple en
prenant pour  X un revêtement purement inséparable de P.   Cepen-%
dant 
  si  k est \uline{de caractéristique nulle} ,  2.12. à
2.1\hand{7}. sont valables sans supposer f une immersion.  En effet, il suffit
de le vérifier pour 2.12. , et  cela résulte alors de 2.11. et du fait
que  pour un préschéma algébrique en c ar 0,  lisse = régulier.



2.3.   \uline{Section hyperplane générique}: \uline{irréductibilité et connexité géomét.}
\uline{Théorème} 3.1. (Bertini-Zariski).  Supposons dim $f(X) \hand{\geq} 2$; et que X
soit  géométriquement irréductible. Alors
la section hyperplane générique $Y_{\hand{\unsure{y|\eta}}}$ a la même propriété.



\mspage{15}
